% =====================================================================
%  CHAPTER 3: 影响生长发育的因素（对应大纲第二章）
%  参考：往届笔记第六章 生长发育影响因素
% =====================================================================
\chapter{影响生长发育的因素}
\setcounter{chapter}{3}
\chapterintro{
\textbf{【双横线/掌握】}遗传因素对生长发育的影响（家族性遗传影响、遗传度、家族聚集性、表观遗传）。\textbf{【难点*】}环境因素（自然地理环境因素、化学性环境污染因素、环境内分泌干扰物、物理性环境污染因素）。\textbf{【单横线/熟悉】}社会决定因素（社会经济和医疗卫生保健、家庭生活质量、文化和教育、现代媒体和娱乐方式）。\textbf{【单横线/熟悉】}行为生活方式因素（生活方式与行为的概念、饮食、运动、睡眠等）。\textbf{【单横线/熟悉】}生长发育长期变化的概念、主要表现、变化原因及其对人类的影响。
}

% ----- 3.1 遗传影响因素 -----
\section{生长发育遗传因素}

\begin{defbox}
\textbf{遗传（heredity）：}形态结构、生理功能及身心发育等性状在亲代与子代间的相似性和连续性。
\end{defbox}

\subsection{家族性遗传影响}

\begin{keybox}
\textbf{一、家族性遗传影响}

1. 生长突增模式、性成熟早晚及月经初潮年龄等与家族遗传密切相关。成年身高与父母平均身高间的遗传度为0.75，即身高75\%取决于遗传因素。

$\triangle$ \textbf{生长发育的家族聚集性（familial aggregation）：}父母与子女身高的相关系数呈随年龄上升的趋势，提示遗传因素在个体越接近成熟的阶段表现得越充分的现象。

2. 遗传对心理行为发展的影响在不同年龄段呈现的表现存在差异。研究发现遗传对感知觉、气质有较直接的影响；而在个性品质、道德行为方面，遗传因素对心理-行为的影响作用随年龄增大而减弱，尤其在青少年阶段，遗传因素的作用远不如环境、教育等因素的影响显著。但环境对某种心理特性或行为的发展所起的作用，往往有赖于这种特性或行为的遗传基础。
\end{keybox}

\subsection{表观遗传}

\begin{keybox}
\textbf{二、表观遗传（epigenetics）：}主要研究不涉及基因核苷酸序列改变的基因表达和调控的可遗传修饰。其遗传方式有DNA序列不变而表型改变，改变有可遗传性和可逆性，不遵循孟德尔遗传定律等特点。表观遗传是环境因素和细胞内的遗传物质间发生交互作用的结果，使生物在环境不断变化的情况下，能维持遗传序列的稳定。
\end{keybox}

% ----- 3.2 物质环境因素 -----
\section{生长发育物质环境因素\textbf{【难点*】}}

\subsection{自然地理环境因素}

\begin{keybox}
\textbf{一、自然地理环境因素：}1. 气候因素；2. 季节因素：春季身高增长快，秋季体重增长快。
\end{keybox}

\subsection{环境污染因素}

\begin{keybox}
\textbf{二、环境污染因素}

\textbf{1. 化学性环境污染因素：}（1）空气污染；（2）铅污染
（3）\textbf{环境内分泌干扰物（environmental endocrine disrupters）：}对机体内天然激素的产生、释放、运输、代谢、消除、结合、功能发挥产生干扰作用，从而破坏机体内环境平衡稳定状态的维持，并对个体机能和发育造成影响的外源性环境化学物质。

①分类：模拟/干扰雌激素、干扰睾酮、干扰甲状腺素、干扰其他内分泌功能。

②人体暴露途径：消化道、呼吸道和皮肤接触。

③对儿童青少年身心健康的影响

A. 生殖毒性，生命早期敏感窗口期暴露，影响子代生殖器官发育和精子成熟。

B. 性发育异常，是儿童性早熟的直接病因或促进因素。

C. 神经毒性，引起持久、不可逆的学习能力缺失和行为发育障碍。

D. 抑制免疫反应。

\textbf{2. 物理性环境污染因素：}噪声污染、电磁辐射污染、放射性污染。
\end{keybox}

% ----- 3.3 社会决定因素 -----
\section{生长发育社会决定因素}

\begin{defbox}
\textbf{一、社会经济和医疗卫生保健：}社会经济状况、医疗卫生保障。

\textbf{二、家庭生活质量：}家庭经济状况、父母受教育程度、家庭结构、家庭氛围、亲子关系、教养方式。

\textbf{三、文化和教育}

\textbf{四、现代媒体和娱乐方式：}电视、网络使用、玩具和读物。
\end{defbox}

% ----- 3.4 行为生活方式因素 -----
\section{生长发育行为生活方式因素}

\begin{defbox}
\textbf{一、生活方式（life styles）：}人们在一定社会文化、经济、风俗、家庭等因素影响下形成的一系列日常生活习惯和生活模式。

\textbf{二、行为（behaviors）：}表现为其具体外显，主要包括饮食（早中晚3：4：3）、运动（体育锻炼）、睡眠、娱乐、消费、社会交往等。
\end{defbox}

% ----- 3.5 促进生长发育的措施 -----
\section{促进生长发育的措施}

\begin{defbox}
综合本章所述影响生长发育的各类因素，促进儿童少年生长发育的主要措施包括：

\begin{enumerate}[leftmargin=*]
    \item \imp{充分利用遗传潜力}——通过遗传咨询和健康教育，帮助家庭了解遗传因素对生长发育的影响，科学认识遗传潜力的发挥
    \item \imp{改善物质环境条件}——治理环境污染，减少铅、环境内分泌干扰物等有害物质的暴露，创造有利于生长发育的自然和人工环境
    \item \imp{优化社会决定因素}——提高社会经济水平，完善医疗卫生保障体系，改善家庭生活质量，提高父母受教育程度，营造良好的家庭氛围和亲子关系
    \item \imp{培养健康行为生活方式}——合理膳食（早中晚能量比3:4:3）、积极参加体育锻炼、保证充足睡眠、合理使用现代媒体
    \item \imp{加强学校健康教育}——发挥学校在健康促进中的核心作用，将健康融入教育全过程
    \item \imp{关注生长发育长期变化趋势}——定期监测群体生长发育状况，及时发现和应对新出现的健康问题（如肥胖率上升等）
\end{enumerate}
\end{defbox}

% ----- 3.6 生长发育长期变化 -----
\section{生长发育的长期变化}

\begin{defbox}
\textbf{生长发育长期变化（secular trend of growth）/ 生长发育长期趋势：}

指以儿童体格发育、青春期和成年身高为核心表现的长期变化。

\textbf{主要表现：}
\begin{enumerate}[leftmargin=*]
    \item \imp{身高、体重的增长}——各年龄组儿童的平均身高和体重逐代增高
    \item \imp{青春期提前}——月经初潮年龄逐代提前
    \item \imp{成人身高的增加}——一代比一代更高
    \item \imp{身体比例的轻微改变}
\end{enumerate}

\textbf{变化原因：}社会经济条件改善、营养水平提高、疾病控制进步、医疗卫生服务改善等综合因素作用的结果。

\textbf{对人类的影响：}生长发育长期变化是人群健康水平提高的正面指标，但也带来新的健康问题（如肥胖率上升等），需要引起重视并采取相应干预措施。
\end{defbox}

% ----- 本章小结 -----
\begin{keybox}
\textbf{本章小结：}

儿童少年生长发育个体差异的形成原因，主要分成两大类：内在的遗传因素和外在的环境因素。
\begin{itemize}[leftmargin=*]
    \item \imp{遗传因素决定生长发育的潜力}，即决定生长发育的可能性
    \item \imp{各种环境因素则在不同程度上影响该潜力的正常发挥}，决定发育的速度及可能达到的程度，即决定生长发育的现实性
\end{itemize}

遗传因素方面，需掌握家族性遗传影响（遗传度为0.75、家族聚集性）和表观遗传（DNA序列不变而表型改变，可遗传、可逆、不遵循孟德尔遗传定律）。环境因素方面，需重点掌握自然地理环境因素（季节因素：春高秋重）、化学性环境污染因素（空气污染、铅污染、环境内分泌干扰物）和物理性环境污染因素（噪声、电磁辐射、放射性污染），其中环境内分泌干扰物为教学难点。社会决定因素方面，需熟悉社会经济状况、医疗卫生保障、家庭生活质量（六方面）、文化和教育、现代媒体和娱乐方式对生长发育的影响。行为生活方式因素方面，需熟悉生活方式与行为的概念及饮食（3:4:3）、运动、睡眠等对生长发育的影响。此外，需熟悉生长发育长期变化的概念、主要表现（身高体重增长、青春期提前、成人身高增加、身体比例改变）、变化原因（社会经济改善、营养提高、疾病控制、医疗卫生进步）及其对人类的影响（正面指标，也带来肥胖等新问题）。
\end{keybox}

% ----- 复习题 -----
\section{复习题}

\begin{enumerate}[leftmargin=*, itemsep=8pt]
    \item \textbf{【名词解释】}遗传（heredity）；家族聚集性；表观遗传（epigenetics）
    \item \textbf{【名词解释】}环境内分泌干扰物（environmental endocrine disrupters）
    \item \textbf{【名词解释】}生活方式（life styles）；生长发育长期变化
    \item \textbf{【简答题】}简述家族性遗传对生长发育的影响（遗传度、家族聚集性）。
    \item \textbf{【简答题】}简述遗传对心理行为发展的影响及其在不同年龄段的表现差异。
    \item \textbf{【简答题】}简述表观遗传的概念、特点及其与生长发育的关系。
    \item \textbf{【简答题/论述题】}试述化学性环境污染因素对儿童少年生长发育的影响（空气污染、铅污染、环境内分泌干扰物）。
    \item \textbf{【简答题】}简述环境内分泌干扰物（EEDs）的分类、暴露途径及对儿童青少年身心健康的影响。
    \item \textbf{【简答题】}简述物理性环境污染因素对儿童少年生长发育的影响。
    \item \textbf{【简答题】}简述社会决定因素对儿童少年生长发育的影响（社会经济和医疗卫生保健、家庭生活质量、文化和教育、现代媒体和娱乐方式）。
    \item \textbf{【简答题】}简述家庭生活质量对儿童生长发育影响的六个方面。
    \item \textbf{【简答题】}简述儿童少年行为生活方式的主要构成及其对健康的影响。
    \item \textbf{【简答题】}试述促进儿童少年生长发育的主要措施。
    \item \textbf{【简答题】}简述生长发育长期变化的概念、主要表现、变化原因及其对人类的影响。
\end{enumerate}

\subsection*{参考答案要点}

\begin{enumerate}[leftmargin=*, itemsep=5pt]
    \item \textbf{遗传}：形态结构、生理功能及身心发育等性状在亲代与子代间的相似性和连续性。\textbf{家族聚集性}：父母与子女身高的相关系数呈随年龄上升的趋势，提示遗传因素在个体越接近成熟的阶段表现得越充分的现象。\textbf{表观遗传}：主要研究不涉及基因核苷酸序列改变的基因表达和调控的可遗传修饰，具有DNA序列不变而表型改变、可遗传性和可逆性、不遵循孟德尔遗传定律等特点。

    \item \textbf{EEDs}：对机体内天然激素的产生、释放、运输、代谢、消除、结合、功能发挥产生干扰作用，从而破坏机体内环境平衡稳定状态的维持，并对个体机能和发育造成影响的外源性环境化学物质。

    \item \textbf{生活方式}：人们在一定社会文化、经济、风俗、家庭等因素影响下形成的一系列日常生活习惯和生活模式。\textbf{生长发育长期变化}：以儿童体格发育、青春期和成年身高为核心表现的长期变化趋势。

    \item 生长突增模式、性成熟早晚及月经初潮年龄等与家族遗传密切相关。成年身高与父母平均身高间的遗传度为0.75，即身高75\%取决于遗传因素。家族聚集性表现为父母与子女身高的相关系数呈随年龄上升的趋势，提示遗传因素在个体越接近成熟的阶段表现得越充分。

    \item 遗传对心理行为发展的影响在不同年龄段呈现的表现存在差异：遗传对感知觉、气质有较直接的影响；而在个性品质、道德行为方面，遗传因素对心理-行为的影响作用随年龄增大而减弱，尤其在青少年阶段，遗传因素的作用远不如环境、教育等因素的影响显著。但环境对某种心理特性或行为的发展所起的作用，往往有赖于这种特性或行为的遗传基础。

    \item 表观遗传主要研究不涉及基因核苷酸序列改变的基因表达和调控的可遗传修饰。特点：DNA序列不变而表型改变，改变具有可遗传性和可逆性，不遵循孟德尔遗传定律。表观遗传是环境因素和细胞内的遗传物质间发生交互作用的结果，使生物在环境不断变化的情况下，能维持遗传序列的稳定。

    \item （1）空气污染：包括室外污染（颗粒物、二氧化氮等）和室内污染（环境烟草暴露、甲醛等），对儿童肺功能和神经系统发育影响广泛；（2）铅污染：铅是环境污染物中毒性最大的重金属之一，任何铅暴露都可能对儿童的健康产生不良影响；（3）EEDs：干扰机体内天然激素的产生、释放、运输、代谢、消除、结合、功能发挥，导致生殖毒性、性发育异常（儿童性早熟的直接病因或促进因素）、神经毒性（持久不可逆的学习能力缺失和行为发育障碍）、抑制免疫反应。

    \item 分类：模拟/干扰雌激素、干扰睾酮、干扰甲状腺素、干扰其他内分泌功能。暴露途径：消化道、呼吸道和皮肤接触。影响：A.生殖毒性——生命早期敏感窗口期暴露，影响子代生殖器官发育和精子成熟；B.性发育异常——是儿童性早熟的直接病因或促进因素；C.神经毒性——引起持久、不可逆的学习能力缺失和行为发育障碍；D.抑制免疫反应。

    \item 三类物理性环境污染因素：（1）噪声污染——损害听觉、影响视觉功能和神经心理行为发育；（2）电磁辐射污染——长期暴露对心血管、消化、内分泌等系统产生功能紊乱甚至损伤器官，低强度慢性射频辐射影响神经系统发育；（3）放射性污染——需达到一定剂量才会发生有害作用。

    \item （1）社会经济和医疗卫生保健：社会经济状况可独立于自然环境因素对生长发育产生直接影响；医疗卫生保障影响卫生资源配置及服务的可及性和公平性；（2）家庭生活质量：家庭经济状况、父母受教育程度、家庭结构、家庭氛围、亲子关系、教养方式；（3）文化和教育：家庭教育、学校教育、社会教育三者共同作用；（4）现代媒体和娱乐方式：电视、网络使用、玩具和读物对认知和行为发育产生影响。

    \item 家庭生活质量的六个方面：（1）家庭经济状况——影响居住环境、膳食营养水平、社交活动和智力投资；（2）父母受教育程度——母亲的教育程度对儿童影响尤为显著；（3）家庭结构——核心家庭和大家庭的儿童优于单亲家庭儿童；（4）家庭氛围——良好温馨的家庭氛围促进健全人格形成，紧张氛围阻碍心理发展；（5）亲子关系——儿童最早建立的人际关系，积极关系促进认知发展；（6）教养方式——情感支持vs过度干预/惩罚对个性发展影响迥异。

    \item 行为是生活方式的具体外显，主要包括：（1）饮食——早中晚能量比3:4:3，合理膳食；（2）运动——体育锻炼促进体格和心理健康；（3）睡眠——消除疲劳、促进体格和神经行为发育；（4）娱乐、消费、社会交往等。健康的行为生活方式是促进儿童少年生长发育的重要保障。

    \item （1）充分利用遗传潜力——科学认识遗传潜力的发挥；（2）改善物质环境条件——治理污染，减少有害物质暴露；（3）优化社会决定因素——提高社会经济水平，改善家庭生活质量；（4）培养健康行为生活方式——合理膳食、体育锻炼、充足睡眠；（5）加强学校健康教育——发挥学校的核心作用；（6）关注生长发育长期变化趋势——定期监测，及时发现和应对新问题。

    \item 概念：以儿童体格发育、青春期和成年身高为核心表现的长期变化趋势。主要表现：（1）身高体重的增长——各年龄组逐代增高；（2）青春期提前——月经初潮年龄逐代提前；（3）成人身高的增加——一代比一代更高；（4）身体比例的轻微改变。变化原因：社会经济条件改善、营养水平提高、疾病控制进步、医疗卫生服务改善等综合因素。对人类的影响：是人群健康水平提高的正面指标，但也带来新问题（如肥胖率上升等），需要引起重视并采取干预措施。
\end{enumerate}

\newpage
